\section{Post Launch Reflections}
\label{sec:PostLaunchReflections}
\subsection{EPQ Example exploration}
\subsubsection*{EPQ Story 1: Raina}

% \textbf{Video Link:} \url{https://web.microsoftstream.com/video/9f5608df-2f6b-4f36-94db-b4adb6031f9f}


\begin{tabularx}{\textwidth}{@{}lX@{}}
    \toprule
    \textbf{Project Title}    & Making clothes using Kitenge, researching African Fashion                                                                      \\
    \midrule
    \textbf{Idea Origin}      & Took inspiration from her love of DT and Kenyan heritage.                                                                      \\
    \textbf{Early Highlights} & Upcycled bucket hat designs shared with friends and family.                                                                    \\
    \textbf{Key Takeaways}    &
    1. Succesfully plan and bring to fruition pices of clothing

    2. Learn to use and repair sewing machines

    3. Learn to validate information sources                                                                                                                   \\
    \textbf{Advice}           &

    1. Don’t stress about choosing the “right” project

    2. If you think you’re done, you're not

    3. Ask questions                                                                                                                                           \\

    \textbf{Final Thoughts}   & Choose a project you’re proud of. Don’t let parents influence your ideas. It’s about skills gained, not future academic links. \\
    \bottomrule
\end{tabularx}

\subsubsection*{EPQ Story 2: Sam}
\begin{tabularx}{\textwidth}{@{}lX@{}}
    \toprule
    \textbf{Project Title}                        & Livestream Esports tournament                                                              \\
    \midrule
    \textbf{Idea Origin}                          & He went through many ideas, always liked games, thought it would be fun to explore further
    \\
    \textbf{Early Highlights}                     & That his idea was feasible
    \\
    \textbf{What if you choose something boring?} &
    It will feel like a waste of time and will likely be
    \\
    \textbf{Advice}                               &
    listen to your mentors advice and start your project early
    \\

    \textbf{Final Thoughts}                       & His EPQ and the skills it gave him were a highlight in his education                       \\
    \bottomrule
\end{tabularx}

\subsubsection*{EPQ Story 3: Lucy}
\begin{tabularx}{\textwidth}{@{}lX@{}}
    \toprule
    \textbf{Project Title}               & A draft of a musical called "Sixth Form: A Student Musical                   \\
    \midrule
    \textbf{Comparison to A levels}      & It was different. Very online and self study, working with peers for advice. \\
    \textbf{Secret to Success>}          & Work on a project you will enjoy                                             \\
    \textbf{What to do when uninspired?} &
    Would tackle a small part of her project which aligned with her hobbies                                             \\
    \textbf{Proud of:}                   &
    Her entire project and her time management skills
    \\

    \textbf{Final Thoughts}              & Keep on top of your EPQ, "a little bit often"
    \\
    \bottomrule
\end{tabularx}


\subsubsection*{EPQ Story 4: Aminah and Michael}
\begin{tabularx}{\textwidth}{@{}lX@{}}
    \toprule
    \textbf{Project Title}    & Spoken Word Poetry Video and Book of drawings and written ideas about Cambridge                         \\
    \midrule
    \textbf{Idea Origin}      &
    Aminah – always liked poetry

    Michael – enjoyed essay subjects and being outside                                                                                  \\
    \textbf{Time Management}  &
    Aminah – made a plan and reflected on feedback

    Michael - find people to work with and bounce ideas off                                                                             \\
    \textbf{Advice}           &
    Try and choose a project you'll enjoy, plan your time, don’t be afraid to change your project and tweak over time
    \\
    \textbf{Lightbulb moment} &
    Peermind maps, breaking down ideas and focusing on elements
    \\

    \textbf{Early Learnings}  & Learnt to edit on Adobe Premiere Rush, learnt different writing styles. Gained social confidence skills \\

    \textbf{words of wisdom}  &
    Aminah – Enjoy your EPQ and plan your time well

    Michael – don’t be afraid to be honest with yourself and don’t worry about changing ideas                                           \\

    \textbf{Independance?}    &
    You don’t need to be independent, ask for help from peers. But also knowing how to work by yourself away from a structured program  \\

    \bottomrule
\end{tabularx}

\subsubsection*{EPQ Story 5: Remy and Jess }
\begin{tabularx}{\textwidth}{@{}lX@{}}
    \toprule

    \textbf{Project Title}          &
    Prompt Based play and Public speaking workshops for primary school pupils

    \\
    \midrule
    \textbf{Working independently?} &
    Remy – difficult at first but enjoyable

    Jess – ups and downs, self-control of project can get dull but branch out into interesting other directions
    \\
    \textbf{Building Ideas}         &
    Remy – enjoyed GCSE drama and plays, built in song lyrics and other hobbies

    Jess – had an idea for project, started building skills to reach her end goal.
    \\
    \textbf{Advice}                 &
    Remy – initial ideas wont be the same as end results. Balance EPQ time with other commitments

    Jess – be flexible and understand its your project and so you can schedule it \\

    \textbf{Proud of:}              &
    Remy – directing skills they have gained and the play itself

    Jess – the networking skills she gained and connections she made              \\
    \bottomrule
\end{tabularx}




\subsubsection*{EPQ Story 6: Sophie }
\begin{tabularx}{\textwidth}{@{}lX@{}}
    \toprule

    \textbf{Project Title}     &
    Is Exposure Therapy The Best Way Of Treating A Specific Phobia?

    \\
    \midrule
    \textbf{Idea Origin?}      &
    Had two ideas and decided on them by looking ahead to university applications

    \\
    \textbf{Dissertation Tips} &
    Take detailed notes and create a research lo. Helps structure your end product

    Focus on the interest in your topic and plan your time well
    \\
    \textbf{Editing Advice}    &
    Asked for family and peer review, took advice and tweaked structure. Don’t be afraid to make large changes
    \\

    \textbf{Uni Benefits:}     &
    The research skills you will gain.
    \\
    \bottomrule
\end{tabularx}


\subsubsection*{EPQ Story 7: Rebecca }
\begin{tabularx}{\textwidth}{@{}lX@{}}
    \toprule

    \textbf{Project Title}     &
    A quiet book

    \\
    \midrule
    \textbf{Idea Origin?}      &
  Wanted a hand made project, integrated child development, enjoyed sowing

    \\
    \textbf{Enjoyed} &
the success of finalised designs
    \\
    \textbf{Evidencing skill development}    &
Spider diagram marked with different skills all rated, during final conclusions remade to demonstrate skills gained / improves
    \\

    \textbf{Uni Benefits:}     &
    Research, independent working skills
    \\

     \textbf{Advice:}     &
    Enjoy the project and you'll be fine
    \\

     \textbf{Mark achieved}     &
    100%
    \\
    \bottomrule
\end{tabularx}



\subsubsection*{EPQ Story 8: Avigayil }
\begin{tabularx}{\textwidth}{@{}lX@{}}
    \toprule

    \textbf{Project Title}     &
    Design and build and autowah guitar pedal

    \\
    \midrule
    \textbf{Idea Origin?}      &
Save cost on a tool she wanted, loves funk music
    \\
    \textbf{Enjoyed} &
The final product and the skills she developed

    \\
    \textbf{Managing workload:}    &
 Broke the construction into stages
    \\

    \textbf{Balancing Time}     &
Used powers hours to give yourself scheduled blocks

    \\

     \textbf{Biggest Challenge:}     &
  Didn't appreciate the full difficulty of skills required

    \\

     \textbf{Favourite Achievement}     &
    The pedal, problem solving and patience skills
    \\
    \bottomrule
\end{tabularx}

\subsection{Artifact Vs Dissertation comparison}
I then watched a couple of video presentations about the variety of end products that can be produced from an EPQ, The findings below have helped me to lean towards a Dissertation style project, as I feel it will allow me to explore my interests in more depth and develop my research skills. 

\begin{tabularx}{\textwidth}{@{}lXX@{}}
\toprule
\textbf{Type of Outcome} & \textbf{Artefact} & \textbf{Dissertation} \\
\midrule



\textbf{Inspiration Taken} &
Building upon existing products &
Taking influence from your own life \\

\textbf{Planning Used} &
Mind maps, Pinterest boards, OneNote to markup scores &
Mind maps, lists, timelines, to-do lists \\

\textbf{Research Contribution} &
Helped understand the scope of the project &
Constantly changing and improving her written drafts \\

\textbf{Skill Progression} &
Gained new skills (e.g., animation); improved time management &
Improved research skills \\
\bottomrule
\end{tabularx}